\documentclass[../../main.tex]{subfiles}% 注意这里的文件路径不能用 ./main.tex ,否则用latexmk编译子文件会报错
\graphicspath{{\subfix{./image/}}{\subfix{../../image/}}} % 指定图片引用路径,后续可以直接使用图片文件名
% 如果图片存放在上级目录,则使用 ../../image/ 作为备用路径

% 例如:
% \begin{figure}[H]
% \centering
% \includegraphics[width=0.3\textwidth]{图.png}
% \caption{}
% \label{figure:图}
% \end{figure}
% 注意:上述\label{}一定要放在\caption{}之后,否则引用图片序号会只会显示??.

\begin{document}

\section{Jordan标准型的进一步讨论和应用}

\begin{theorem}\label{theorem:关于一个特征值的Jordan块的个数等于其度数,阶数直和等于其重数}
线性变换 $\varphi$ 的特征值 $\lambda_1$ 的几何重数(度数)等于 $\varphi$ 的 Jordan 标准型中属于特征值 $\lambda_1$ 的 Jordan 块的个数，$\lambda_1$ 的代数重数等于所有属于特征值 $\lambda_1$ 的 Jordan 块的阶数之和。并且属于特征值$\lambda_1$的每一个Jordan块都有且仅有一个线性无关的特征向量.
\end{theorem}
\begin{proof}
设 $V$ 是 $n$ 维复线性空间，$\varphi$ 是 $V$ 上的线性变换。设 $\varphi$ 的初等因子组为
\begin{align}
(\lambda - \lambda_1)^{r_1}, (\lambda - \lambda_2)^{r_2}, \cdots, (\lambda - \lambda_k)^{r_k}, \label{equation--7.7.1}
\end{align}
\refthe{theorem:线性变换一定有一个Jordan矩阵作为其表示矩阵}告诉我们，存在 $V$ 的一组基 $\{e_{11}, e_{12}, \cdots, e_{1r_1}; e_{21}, e_{22}, \cdots, e_{2r_2}; \cdots; e_{k1}, e_{k2}, \cdots, e_{kr_k}\}$，使得 $\varphi$ 在这组基下的表示矩阵为
\begin{align*}
\boldsymbol{J} = 
\begin{pmatrix}
\boldsymbol{J}_1 & & & \\
& \boldsymbol{J}_2 & & \\
& & \ddots & \\
& & & \boldsymbol{J}_k
\end{pmatrix}.
\end{align*}
上式中每个 $\boldsymbol{J}_i$ 是相应于初等因子 $(\lambda - \lambda_i)^{r_i}$ 的 Jordan 块，其阶正好为 $r_i$。令 $V_i$ 是由基向量 $e_{i1}, e_{i2}, \cdots, e_{ir_i}$ 生成的子空间，则
\begin{gather}\label{equation--7.7.2}
\begin{aligned}
\varphi(e_{i1}) &= \lambda_i e_{i1}, 
 \\
\varphi(e_{i2}) &= e_{i1} + \lambda_i e_{i2},  \\
&\cdots\cdots\cdots\cdots \\
\varphi(e_{ir_i}) &= e_{i,r_i - 1} + \lambda_i e_{ir_i}.
\end{aligned}
\end{gather}
这表明 $\varphi(V_i) \subseteq V_i$，即 $V_i (i = 1, 2, \cdots, k)$ 是 $\varphi$ 的不变子空间。显然我们有
\[
V = V_1 \oplus V_2 \oplus \cdots \oplus V_k.
\]

线性变换 $\varphi$ 限制在 $V_1$ 上 (仍记为 $\varphi$) 便成为 $V_1$ 上的线性变换。这个线性变换在基 $\{e_{11}, e_{12}, \cdots, e_{1r_1}\}$ 下的表示矩阵为
\begin{align*}
\boldsymbol{J}_1 = 
\begin{pmatrix}
\lambda_1 & 1 & & & \\
& \lambda_1 & 1 & & \\
& & \ddots & \ddots & \\
& & & \ddots & 1 \\
& & & & \lambda_1
\end{pmatrix}.
\end{align*}
注意到 $\boldsymbol{J}_1$ 的特征值全为 $\lambda_1$，并且 $\lambda_1\boldsymbol{I} - \boldsymbol{J}_1$ 的秩等于 $r_1 - 1$，故 $\boldsymbol{J}_1$ 只有一个线性无关的特征向量，不妨选为 $e_{11}$。显然$e_{11}$(补0)延拓到$n$维得到的向量$\hat{e}_{11}$ 也是 $\varphi$ 作为 $V$ 上线性变换关于特征值 $\lambda_1$ 的特征向量。不失一般性，不妨设在 $\varphi$ 的初等因子组即\eqref{equation--7.7.1} 式中
\[
\lambda_1 = \lambda_2 = \cdots = \lambda_s, \quad \lambda_i \neq \lambda_1 (i = s + 1, \cdots, k),
\]
则 $\boldsymbol{J}_1, \cdots, \boldsymbol{J}_s$ 都以 $\lambda_1$ 为特征值，且{\heiti 相应于每一块有且只有一个线性无关的特征向量}。相应的特征向量可取为
\begin{align}
\hat{e}_{11}, \hat{e}_{21}, \cdots, \hat{e}_{s1}, \label{equation--7.7.3}
\end{align}
其中$\hat{e}_{i1}$都是$e_{i1}$(补0)延拓到$n$维得到的向量.
显然这是 $s$ 个线性无关的特征向量。如果 $\lambda_i \neq \lambda_1$，则容易看出 $\mathrm{r}(\lambda_1\boldsymbol{I} - \boldsymbol{J}_i) = r_i$，于是
\begin{align*}
\mathrm{r}(\lambda_1\boldsymbol{I} - \boldsymbol{J}) &= \sum_{i = 1}^{k} \mathrm{r}(\lambda_1\boldsymbol{I} - \boldsymbol{J}_i) = (r_1 - 1) + \cdots + (r_s - 1) + r_{s + 1} + \cdots + r_k = n - s.
\end{align*}
因此 $\varphi$ 关于特征值 $\lambda_1$ 的特征子空间 $V_{\lambda_1}$ 的维数等于 $n - \mathrm{r}(\lambda_1\boldsymbol{I} - \boldsymbol{J}) = s$，从而特征子空间 $V_{\lambda_1}$ 以\eqref{equation--7.7.3}式中的向量为一组基。又 $\lambda_1$ 是 $\varphi$ 的 $r_1 + r_2 + \cdots + r_s$ 重特征值，因此 $\lambda_1$ 的重数与度数之差等于
\[
(r_1 + r_2 + \cdots + r_s) - s.
\]
\end{proof}

\begin{definition}[根子空间]
设 $\lambda_0$ 是 $n$ 维复线性空间 $V$ 上线性变换 $\varphi$ 的特征值，则
\begin{align*}
R(\lambda_0) = \{\boldsymbol{v} \in V \mid (\varphi - \lambda_0\boldsymbol{I})^n(\boldsymbol{v}) = \boldsymbol{0}\}
\end{align*}
构成了 $V$ 的一个子空间，称为属于特征值 $\lambda_0$ 的\textbf{根子空间}。 
\end{definition}

\begin{lemma}[根子空间维数]\label{lemma:根子空间维数}
设 \( A \in \mathbb{F}^{n \times n} \) 且 \( \mathbb{F} \) 是域. 若 \( \lambda \in \mathbb{F} \) 是 \( A \) 的 \( r \) 重特征值, 则我们有
\begin{align}
\dim \mathrm{Ker} \left( \lambda E - A \right)^m = r, \forall m \geqslant r. \label{eq::::------23.7}
\end{align}
\end{lemma}
\begin{proof}
注意到
\[
\dim \mathrm{Ker} \left( \lambda E - A \right)^r = n - \mathrm{rank} \left( \left( \lambda E - A \right)^r \right).
\]
而秩不随域扩张而改变, 故不妨设 \( \mathbb{F} \) 是代数闭域. 则不妨设 \( A \) 为 Jordan 标准型
\[
\begin{pmatrix}
J_{n_1} \left( \lambda \right) & & & \\
& \ddots & & \\
& & J_{n_s} \left( \lambda \right) & \\
& & & J
\end{pmatrix}, \quad \sum_{i=1}^s n_i = r,
\]
这里 \( J \) 没有特征值 \( \lambda \). 于是直接计算有
\begin{gather*}
\left( A - \lambda E \right)^r = \begin{pmatrix}
J_{n_1}^r \left( 0 \right) & & & \\
& \ddots & & \\
& & J_{n_s}^r \left( 0 \right) & \\
& & & \left( J - \lambda E_{n-r} \right)^r
\end{pmatrix} = \begin{pmatrix}
0 & & & \\
& \ddots & & \\
& & 0 & \\
& & & \left( J - \lambda E_{n-r} \right)^r
\end{pmatrix}, 
\\
\det \left( J - \lambda E_{n-r} \right)^r \neq 0.
\end{gather*}
现在知 \( \mathrm{rank} \left( A - \lambda E \right)^r = n - r \), 这就得到了 \eqref{eq::::------23.7} 的 \( m = r \) 的情况. 注意到当 \( m > r \) 时, 上述矩阵 \( 0 \) 块在次方时不会继续丢失阶数,而$\left( J - \lambda E \right)^m$仍可逆, 因此秩不会再次减少, 所以 \eqref{eq::::------23.7} 对任何 \( m \geqslant r \) 都成立, 这就完成了证明.
\end{proof}

\begin{theorem}\label{theorem:Jordan标准型的几何意义}
设 $\varphi$ 是 $n$ 维复线性空间 $V$ 上的线性变换。
\begin{enumerate}[(1)]
\item 若 $\varphi$ 的初等因子组为
\[
(\lambda - \lambda_1)^{r_1}, (\lambda - \lambda_2)^{r_2}, \cdots, (\lambda - \lambda_k)^{r_k},
\]
则 $V$ 可分解为 $k$ 个$\varphi$的不变子空间的直和：
\begin{align}
V = V_1 \oplus V_2 \oplus \cdots \oplus V_k, \label{equation--7.7.5}
\end{align}
其中 $V_i$ 是维数等于 $r_i$ 的关于 $\varphi - \lambda_i\boldsymbol{I}$ 的循环子空间；

\item 若 $\lambda_1, \cdots, \lambda_s$ 是 $\varphi$ 的全体不同特征值，则 $V$ 可分解为 $s$ 个不变子空间的直和：
\begin{align*}
V = R(\lambda_1) \oplus R(\lambda_2) \oplus \cdots \oplus R(\lambda_s), 
\end{align*}
其中 $R(\lambda_i)$ 是 $\lambda_i$ 的根子空间，$R(\lambda_i)$ 的维数等于 $\lambda_i$ 的代数重数，且每个 $R(\lambda_i)$ 又可分解为\eqref{equation--7.7.5} 式中若干个 $V_j$ 的直和。 
\end{enumerate} 
\end{theorem}
\begin{proof}
在\refthe{theorem:关于一个特征值的Jordan块的个数等于其度数,阶数直和等于其重数}的证明的基础上,现在再来看 $\boldsymbol{J}_1$ 所对应的子空间 $V_1$，由 \eqref{equation--7.7.2}式中诸等式可知
\[
(\varphi - \lambda_1\boldsymbol{I})(e_{1r_1}) = e_{1,r_1 - 1}, \cdots, (\varphi - \lambda_1\boldsymbol{I})(e_{12}) = e_{11}, (\varphi - \lambda_1\boldsymbol{I})(e_{11}) = \boldsymbol{0},
\]
因此，若记 $\boldsymbol{\alpha} = e_{1r_1}, \psi = \varphi - \lambda_1\boldsymbol{I}$，则
\begin{align}\label{eq:24gf34t4g3423gr6745}
\psi(\boldsymbol{\alpha}) = e_{1,r_1 - 1},  \psi^2(\boldsymbol{\alpha}) = e_{1,r_1 - 2},  \cdots,  \psi^{r_1 - 1}(\boldsymbol{\alpha}) = e_{11},  \psi^{r_1}(\boldsymbol{\alpha}) = \boldsymbol{0}.
\end{align}
也就是说
\[
\{\boldsymbol{\alpha}, \psi(\boldsymbol{\alpha}), \psi^2(\boldsymbol{\alpha}), \cdots, \psi^{r_1 - 1}(\boldsymbol{\alpha})\}
\]
构成了 $V_1$ 的一组基。

上面的事实说明，每个 Jordan 块 $\boldsymbol{J}_i$ 对应的子空间 $V_i$ 是一个循环子空间。把属于同一个特征值，比如属于 $\lambda_1$ 的所有循环子空间加起来构成 $V$ 的一个子空间：
\[
R(\lambda_1) = V_1 \oplus \cdots \oplus V_s.
\]
若 $\boldsymbol{v} \in R(\lambda_1)$，则由\eqref{eq:24gf34t4g3423gr6745}式可知 $(\varphi - \lambda_1\boldsymbol{I})^s(\boldsymbol{v}) = \boldsymbol{0}$，其中
\[
s = \dim R(\lambda_1) = r_1 + \cdots + r_s.
\]
事实上，我们可以证明
\begin{align}\label{equation--7.7.4}
R(\lambda_1) = \{\boldsymbol{v} \in V \mid (\varphi - \lambda_1\boldsymbol{I})^n(\boldsymbol{v}) = \boldsymbol{0}\}.
\end{align}
为证明 \eqref{equation--7.7.4}式成立，设 $U = \{\boldsymbol{v} \in V \mid (\varphi - \lambda_1\boldsymbol{I})^n(\boldsymbol{v}) = \boldsymbol{0}\}$，则由上面的分析知道，$R(\lambda_1) \subseteq U$。另一方面，任取 $\boldsymbol{v} \in U$，设 $\boldsymbol{v} = \boldsymbol{v}_1 + \boldsymbol{v}_2$，其中 $\boldsymbol{v}_1 \in R(\lambda_1)$，$\boldsymbol{v}_2 \in V_{s + 1} \oplus \cdots \oplus V_k$。因为 $(\lambda - \lambda_1)^n$ 与 $(\lambda - \lambda_{s + 1})^n \cdots (\lambda - \lambda_k)^n$ 互素，故存在多项式 $p(\lambda), q(\lambda)$，使
\[
(\lambda - \lambda_1)^n p(\lambda) + (\lambda - \lambda_{s + 1})^n \cdots (\lambda - \lambda_k)^n q(\lambda) = 1.
\]
将 $\lambda = \varphi$ 代入上式并作用在 $\boldsymbol{v}$ 上可得
\begin{align*}
\boldsymbol{v} &= p(\varphi)(\varphi - \lambda_1\boldsymbol{I})^n(\boldsymbol{v}) + q(\varphi)(\varphi - \lambda_{s + 1}\boldsymbol{I})^n \cdots (\varphi - \lambda_k\boldsymbol{I})^n(\boldsymbol{v}) \\
&= q(\varphi)(\varphi - \lambda_{s + 1}\boldsymbol{I})^n \cdots (\varphi - \lambda_k\boldsymbol{I})^n(\boldsymbol{v}_1) + q(\varphi)(\varphi - \lambda_{s + 1}\boldsymbol{I})^n \cdots (\varphi - \lambda_k\boldsymbol{I})^n(\boldsymbol{v}_2) \\
&= q(\varphi)(\varphi - \lambda_{s + 1}\boldsymbol{I})^n \cdots (\varphi - \lambda_k\boldsymbol{I})^n(\boldsymbol{v}_1) \in R(\lambda_1).
\end{align*}
这就证明了 \eqref{equation--7.7.4}式。

上面的结果表明：特征值 $\lambda_0$ 的根子空间可表示为若干个循环子空间的直和，每个循环子空间对应于一个 Jordan 块。
虽然我们前面的讨论是对特征值 $\lambda_1$ 进行的，其实对任一特征值 $\lambda_i$ 均适用。 
\end{proof}

\begin{proposition}
证明：复数域上的方阵 $\boldsymbol{A}$ 必可分解为两个对称阵的乘积。
\end{proposition}
\begin{proof}
设 $\boldsymbol{P}$ 是非异阵且使 $\boldsymbol{P}^{-1}\boldsymbol{AP} = \boldsymbol{J}$ 为 $\boldsymbol{A}$ 的 Jordan 标准型，于是 $\boldsymbol{A} = \boldsymbol{PJP}^{-1}$。设 $\boldsymbol{J}_i$ 是 $\boldsymbol{J}$ 的第 $i$ 个 Jordan 块，则
\begin{align*}
\boldsymbol{J}_i=\left( \begin{matrix}
\lambda _i&		1&		&		&		\\
&		\lambda _i&		1&		&		\\
&		&		\ddots&		\ddots&		\\
&		&		&		\ddots&		1\\
&		&		&		&		\lambda _i\\
\end{matrix} \right) =\left( \begin{matrix}
&		&		&		1&		\lambda _i\\
&		&		1&		\lambda _i&		\\
&		\begin{turn}{80}$\ddots$\end{turn}&		\begin{turn}{80}$\ddots$\end{turn}&		&		\\
1&		\begin{turn}{80}$\ddots$\end{turn}&		&		&		\\
\lambda _i&		&		&		&		\\
\end{matrix} \right) \left( \begin{matrix}
&		&		&		&		1\\
&		&		&		1&		\\
&		&		\begin{turn}{80}$\ddots$\end{turn}&		&		\\
&		1&		&		&		\\
1&		&		&		&		\\
\end{matrix} \right) ,
\end{align*}
即 $\boldsymbol{J}_i$ 可分解为两个对称阵之积。因此 $\boldsymbol{J}$ 也可以分解为两个对称阵之积，记为 $\boldsymbol{S}_1, \boldsymbol{S}_2$，于是
\[
\boldsymbol{A} = \boldsymbol{PJP}^{-1} = \boldsymbol{PS}_1\boldsymbol{S}_2\boldsymbol{P}^{-1} = (\boldsymbol{PS}_1\boldsymbol{P}')(\boldsymbol{P}^{-1})'\boldsymbol{S}_2\boldsymbol{P}^{-1}. 
\] 
显然$\boldsymbol{PS}_1\boldsymbol{P}'$和$(\boldsymbol{P}^{-1})'\boldsymbol{S}_2\boldsymbol{P}^{-1}$都是对称矩阵,
故$\boldsymbol{A}$ 必可分解为两个对称阵的乘积。
\end{proof}

\begin{example}
已知
\[
\boldsymbol{A} = 
\begin{pmatrix}
3 & 1 & 0 & 0 \\
-4 & -1 & 0 & 0 \\
6 & 1 & 2 & 1 \\
-14 & -5 & -1 & 0
\end{pmatrix},
\] 
计算$A^k.$
\end{example}
\begin{solution}
用初等变换把 $\lambda\boldsymbol{I} - \boldsymbol{A}$ 化为对角 $\lambda$-矩阵并求出它的初等因子组为
\[
(\lambda - 1)^2, (\lambda - 1)^2.
\]
因此，$\boldsymbol{A}$ 的 Jordan 标准型为
\begin{align*}
\boldsymbol{J} = 
\begin{pmatrix}
1 & 1 & & \\
0 & 1 & & \\
 & & 1 & 1 \\
 & & 0 & 1
\end{pmatrix}.
\end{align*} 
因为
\begin{align*}
\boldsymbol{P}^{-1}\boldsymbol{A}^k\boldsymbol{P} = (\boldsymbol{P}^{-1}\boldsymbol{AP})^k = \boldsymbol{J}^k,
\end{align*}
故先计算 $\boldsymbol{J}^k$。注意 $\boldsymbol{J}$ 是分块对角阵，它的 $k$ 次方等于将各对角块 $k$ 次方，因此
\begin{align*}
\boldsymbol{J}^k = 
\begin{pmatrix}
1 & 1 & & \\
0 & 1 & & \\
    & & 1 & 1 \\
    & & 0 & 1
\end{pmatrix}^k
=
\begin{pmatrix}
1 & k & & \\
0 & 1 & & \\
    & & 1 & k \\
    & & 0 & 1
\end{pmatrix},
\end{align*}
\begin{align*}
\boldsymbol{A}^k = \boldsymbol{PJ}^k\boldsymbol{P}^{-1} 
&=
\begin{pmatrix}
1 & 0 & 0 & 0 \\
-2 & 1 & 0 & 0 \\
1 & 0 & 1 & 0 \\
-5 & 0 & -1 & 1
\end{pmatrix}
\begin{pmatrix}
1 & k & & \\
0 & 1 & & \\
    & & 1 & k \\
    & & 0 & 1
\end{pmatrix}
\begin{pmatrix}
1 & 0 & 0 & 0 \\
-2 & 1 & 0 & 0 \\
1 & 0 & 1 & 0 \\
-5 & 0 & -1 & 1
\end{pmatrix}^{-1} \\
&=
\begin{pmatrix}
2k + 1 & k & 0 & 0 \\
-4k & -2k + 1 & 0 & 0 \\
6k & k & k + 1 & k \\
-14k & -5k & -k & -k + 1
\end{pmatrix}.
\end{align*}

\end{solution}


























































































































\end{document}